Esta memoria recoge las actividades desarrolladas durante las practicas curriculares realizadas en ActLab dentro del proyecto \textit{Didascalia VC - ACTLAB}, una aplicacion de Realidad Virtual orientada al entrenamiento y simulacion de contextos educativos.

El trabajo se ha centrado en cinco ejes tecnicos principales: optimizacion inicial de rendimiento grafico en VR para reducir tirones y estabilizar FPS, diseno de una herramienta basada en grafos para describir comportamientos del alumnado, refactorizacion y robustecimiento del sistema de conexion con la web y su procesamiento de mensajes, ampliacion y validacion funcional del sistema de conflictos (incluyendo nuevos casos TEA y TDAH), y despliegue/verificacion multiplataforma con objetivo Meta Quest 3.

Adicionalmente, se han integrado mejoras de flujo de trabajo para el equipo, como la seleccion de servidores de prueba (entornos de desarrollo y produccion), y se han realizado tareas de integracion de contenido (animaciones y adaptacion de prefabs) necesarias para consolidar los nuevos comportamientos incorporados.

La memoria documenta tanto los resultados tecnicos obtenidos como la metodologia seguida, relacionando las tareas realizadas con competencias adquiridas durante el grado y dejando identificados los puntos que requieren completado administrativo antes de su entrega final en GIPE.